\pagenumbering{gobble}
\pagestyle{emptyusfq}

\begin{center}
{\bfseries\fontsize{18}{22}\selectfont GUÍA Y PLANTILLA PARA ENTREGA DE TRABAJO FINAL DE POSGRADO DE LA USFQ}
\end{center}

\USFQNotaPlantilla{Esta es una hoja de instrucciones generales que debe ser borrada antes de entregar el trabajo. Para ocultarla en LaTeX, cambie \texttt{\textbackslash USFQIncluirGuiatrue} por \texttt{\textbackslash USFQIncluirGuiafalse} en \texttt{preambulo/preambulo.tex}.}

\begin{enumerate}
  \item Todos los textos en rojo deben ser reemplazados por los datos particulares e información correspondiente y cambiados a negro antes de entregar el trabajo.
  \item Se debe escoger uno de los estilos que se listan a continuación para el formato del trabajo de titulación y referencias bibliográficas, o aquel que se ajuste a las necesidades del área de conocimiento en la que está trabajando.
\end{enumerate}

\USFQNotaPlantilla{Se debe seleccionar un solo estilo y se deberá cumplir de manera estricta todos los parámetros de ese estilo.}

\begin{itemize}
  \item \textbf{APA}, según \textit{Publication manual of the American psychological association}.
  \item \textbf{Chicago}, según \textit{The Chicago Manual of Style}.
  \item \textbf{Citing Medicine}, según \textit{Citing Medicine}.
  \item \textbf{MLA}, según \textit{MLA Style Manual and Guide to Scholarly Publishing}.
  \item \textbf{Oxford}, según \textit{New Hart's Rules: The Oxford Style Guide}.
\end{itemize}

\begin{enumerate}
  \setcounter{enumi}{3}
  \item El esquema es solo referencial, pueden tener secciones y subtítulos adicionales. Se recomienda el siguiente esquema:
\end{enumerate}

\begin{table}[H]
\centering
\caption*{Sugerencia USFQ para niveles de títulos}
\begin{singlespace}
\begin{tblr}{
  width=\textwidth,
  colspec={p{0.18\textwidth} X[j]},
}
Nivel 1 & \textbf{CENTRADO, EN NEGRILLA, TODAS LAS PALABRAS EN MAYÚSCULAS, tamaño de letra 12 puntos} \\
Nivel 2 & \textbf{Alineado a la izquierda, en negrilla, mayúscula solo la primera letra de la primera palabra, tamaño de letra 12 puntos} \\
Nivel 3 & \textbf{Con sangría de párrafo, negrilla, mayúscula solo la primera letra de la primera palabra, tamaño de letra 12 puntos y terminado en punto.} \\
Nivel 4 & \textbf{\textit{Con sangría de párrafo, negrilla, cursiva, mayúscula solo la primera letra de la primera palabra, tamaño de letra 12 puntos y terminado en punto.}} \\
Nivel 5 & \textit{Con sangría de párrafo, cursiva, mayúscula solo la primera letra de la primera palabra, tamaño de letra 12 puntos y terminado en punto.} \\
\end{tblr}
\end{singlespace}
\end{table}

\begin{enumerate}
  \setcounter{enumi}{4}
  \item Todos los textos del trabajo que no sean títulos deben tener un tamaño de letra de 12 puntos.
  \item Utilizar el formato de hoja A4.
  \item Escribir el texto a espacio doble.
  \item Todas las páginas del trabajo de titulación deben ir numeradas en la parte superior derecha, excepto la portada.
\end{enumerate}

\USFQNotaPlantilla{Se espera que todos los trabajos sean entregados siguiendo uno de los formatos establecidos en este documento. Se deberá respetar el formato en relación a tipo y tamaño de fuentes, espaciado, márgenes, etc. Se recomienda que el desarrollo del tema del trabajo de titulación no tenga una extensión mayor a 30.000 palabras.}

\USFQNotaPlantilla{Los estudiantes son los responsables directos de verificar que todo el material que utilizan en su trabajo de titulación, en especial tablas y figuras, puedan ser usados y reproducidos. Material que tenga derechos de autor reservados no deberá ser incluido en este trabajo a menos que el estudiante cuente con el permiso explícito de sus autores, con excepción de citas y paráfrasis, las que deben ser siempre acreditadas a su autor.}

\USFQNotaPlantilla{La redacción debe ser escrita en lenguaje académico y las fuentes deben ser citadas siguiendo el formato escogido, tanto para citas en el texto como para las Referencias Bibliográficas. En esta plantilla se conserva la lógica BibTeX con estilo APA.}

\USFQNotaPlantilla{\textbf{ATENCIÓN:} Si el trabajo de titulación es un producto publicable escrito en los formatos reconocidos internacionalmente para publicación, por ejemplo artículo científico, \textit{journal} o \textit{paper}, no es necesario que sea modificado a este formato. Únicamente es requisito aumentar todas las hojas preliminares presentadas en la plantilla a continuación.}

\begin{center}
\bfseries
A PARTIR DE LA SIGUIENTE HOJA, SE ENCUENTRA UNA PLANTILLA EN LA QUE SE PUEDE REEMPLAZAR LOS TEXTOS EN ROJO CON LA INFORMACIÓN PARA PRESENTAR EL TRABAJO FINAL.
\end{center}

\clearpage
